\documentclass[a4paper]{article}
\usepackage[pdftex]{graphicx}
\usepackage[utf8]{inputenc}
\usepackage{enumerate}
\usepackage{icomma}
\usepackage{amssymb}
\usepackage{eurosym}
\usepackage{tikz}
\usepackage{href-ul}
\hypersetup{
	colorlinks=true,
	linkcolor=blue,
	urlcolor=blue}
\usepackage{geometry}
\geometry{a4paper, top=15mm, left=15mm, right=15mm, bottom=15mm,
	headsep=10mm, footskip=12mm}

\begin{document}
\begin{tikzpicture}(10,3)
	\draw[ultra thick](2,0) --(10,0) -- (10,1.3) --(2,1.3) -- (2,0);
	\draw[fill=black](2,0)-- (0,.6) -- (2,.6) -- (2,0);
	\node at (6,.6) {\large Lösungsblatt von \href{https://www.okuyakl.de}{www.okuyakl.de}};
\end{tikzpicture}
\vspace{0.5 cm}

\noindent{\bf Aufgabe 1.}\\
Wir stellen um und logarithmieren:
$$ 
\renewcommand{\arraystretch}{2}
\begin{array}{rcll} 
{2\over 5} &=& \left({2 \over 1,2^2}\right)^x &| \lg()\\
\lg {2\over 5} &=& x\cdot \lg {2 \over 1,2^2} & | : \lg {2 \over 1,2^2} \\
{\lg {2\over 5}\over \lg {2 \over 1,2^2} }&=& x \\
x &=& -2,8
\end{array}
$$

\noindent{\bf Aufgabe 2.}\\
Es handelt sich um eine exponentielle Abnahme; der Wachstumsfaktor $a$ ist:
$$a=\left( 1 - {1,5\over 100}\right)=0,985$$
Die Zerfallsgleichung lautet damit:
$$N(t)=140 \cdot 0,985^t$$
Für die Aufgabenstellung setzen wir die Anzahl der aktiven Kerne zu Beginn auf 100\% und am Ende
des Beobachtungszeitraums auf 42\% und erhalten:
$$
\renewcommand{\arraystretch}{2}
\begin{array}{rcll} 
42\%&=&100\% \cdot 0,985^t & | : 100\% \\
0,42 &=& 0,985^t &| \lg()\\
\lg 0,42 &=& t \cdot \lg 0,985 &| : \lg 0,985\\
t &=& {\lg 0,42\over \lg 0,985} &= 57,4\\
\end{array}
$$
Antwort: nach $57,4~s$ beträgt die Anzahl aktiver Kerne noch $42\%$ des Anfangswertes.

\noindent{\bf Aufgabe 3. a)}\\
$$\log_5(\sqrt[4]{5ab})= {1 \over 4} \cdot (\log_5 5 + \log_5 a + \log_5 b ) 
= \underline{{1 \over 4} + {1 \over 4} \cdot (\log_5 a + \log_5 b )}$$
\noindent {\bf Aufgabe 3. b)}\\
$$\log_a \left( {14 a \over b^4} \right)=\log_a 14 + \log_a a - 4 \cdot \log b=
\underline{ 1+\log_a 14  - 4 \cdot \log b}$$
\noindent {\bf Aufgabe 4.}\\
$$\dots = -0,25 \log_a b - 0,5 \log_a b +0,4 \log_a b^{-1,5}=
\underline{-1,35 \log_a b}$$
\noindent {\bf Aufgabe 5.}\\
$$
\renewcommand{\arraystretch}{2}
\begin{array}{rcll} 
0,25 \cdot 20000 &=& 20000 \cdot 3^{-0,3t} & | : 20000\\
0,25 &=& 3^{-0,3t} & | \log(\dots)\\
\log 0,25 &=& -0,3t \cdot \log 3 & | : \log 3\\
{\log 0,25 \over \log 3} &=& -0,3t & |: (-0,3)\\ 
{\log 0,25 \over  -0,3 \log 3} &=& t \\
\underline{4,2} &=& t 
\end{array}
$$
Nach etwa 4 Jahren beträgt der Bestand nur noch 25\% der Anfangspopulation.

\noindent {\bf Aufgabe 6.}\\
$$
\renewcommand{\arraystretch}{2}
\begin{array}{rcll} 
5^{2x+1} - 3 \cdot 5^{x+2} + 250 &=& 0 & \\
5 \cdot 5^{2x} - 3\cdot 5^2 \cdot 5^x + 250 &=& 0 & |\textnormal{Substitution:}\quad 5^x=u\\
5 \cdot u^2 -75 u +250 &=& 0 & |:5\\
u^2 -15 u + 50 &=& 0 & |\textnormal{Mitternachtsformel}\\
u_{1/2} &=& \frac{15 \pm \sqrt{225-4\cdot 50}}{2}\\
u_1 &=& 10\\
u_2 &=& 5 &| \textnormal{Rücksubstitution} \quad u=5^x\\
10 &=& 5^{x_1} \quad \Rightarrow \quad \underline{x_1} &= \underline{\log_5 10} \\
5 &=& 5^{x_2} \quad \Rightarrow \quad \underline{x_2} &= \underline{1} \\
\end{array}
$$

\noindent {\bf Aufgabe 7.}\\
Die Ersatzgleichung lautet:
$$
\begin{array}{rcl} 
x^{2 \over 3} = & 0,5^{1 \over 3} &  | ()^3 \\
x^2 = &  0,5 & | \sqrt{\quad} \\
x  = & \sqrt{1 \over 2} & = {\sqrt{2} \over 2}  
\end{array}
$$

\noindent {\bf Aufgabe 8.}\\
Wir wenden die Logarithmengesetze an:
$$
\begin{array}{rl} 
T(x)= & 5 \cdot \log_5{x} + 2x \cdot \log_5{5} + 5 \cdot \log_5{5^2}-5 \cdot \log_5{x}\\
= & 2x \cdot 1 + 5 \cdot 2 = 2x + 10  
\end{array}
$$

\noindent {\bf Aufgabe 9. a)}\\
Aufstellen der Funktionsgleichung für exponentielle Abnahme:
$$ 
\begin{array}{rl}
f(t)=& f(0)\cdot \left( 1 - {p \over 100} \right)^t \\
f(t)=& 50000 \cdot 0,91^t
\end{array}
$$

\noindent {\bf Aufgabe 9. b)}\\
Wenn der Fischbestand um 20000 abnimmt, sind noch 30000 vorhanden:
$$ 
\begin{array}{rcll}
30000&=& 50000 \cdot 0,91^t & | :50000\\
0,6  &=& 0,91^t & | \log{} \\
\log{0,6} &=& t \cdot \log{0,91} &| : \log{0,91}\\
t &=& 5,4 &
\end{array}
$$
Es dauert 5,4 Jahre, bis die Behörden Maßnahmen ergreifen.

\noindent {\bf Aufgabe 9. c)}\\
Exponentielles Wachstum mit Faktor $ \left( 1 + {6 \over 100} \right) $ ab Zeitpunkt $t=5,4$:
$$ 
\begin{array}{rcll}
50000&=& 30000 \cdot 1,06^t & | :30000\\
{5 \over 3}  &=& 1,06^t & | \log{} \\
\log{5 \over 3} &=& t \cdot \log{1,06} &| : \log{1,06}\\
t &=& 8,8 &
\end{array}
$$
Es dauert insgesamt 5,4 + 8,8 = 14,2 Jahre, bis sich der Fischbestand wieder erholt hat.

\noindent {\bf Aufgabe 10.}\\
Die Funktionsgleichung $y=b\cdot a^x$ besteht aus dem Startwert b=\EUR {47900} und dem Abnahmefaktor:
$$a=1-{p \over 100} = 0,77$$
Beides eingesetzt ergibt:
$$y=\textnormal{\EUR{47900}}\cdot 0,77^x$$
Das Auto hat seinen Restwert nach $x$ Jahren: 
$$
\renewcommand{\arraystretch}{1.5}
\begin{array}{rcll}
\textnormal{\EUR{4550}} &=& \textnormal{\EUR{47900}}\cdot 0,77^x & |:\textnormal{\EUR{47900}}\\
0,095 &=& 0,77^x & |\log(\dots)\\
\log 0,095&=& x\cdot \log 0,77 | : \log 0,77 \\
{\log 0,095 \over \log 0,77}&=& x \\
x &=& \underline{9}  
\end{array}
$$
Das Fahrzeug ist 9 Jahre alt.

\noindent {\bf Aufgabe 11.}\\
Das Kapital wächst exponentiell bis zur Ver-n-fachung mit folgender Gleichung:
$$
\renewcommand{\arraystretch}{1.5}
\begin{array}{rcll}
n \cdot K_0 &=& K_0 \cdot \left(1+{p \over 100}\right)^x & | : K_0 \\
n &=& \left(1+{p \over 100}\right)^x & | \log(\dots) \\
\log n &=& x \cdot \log \left(1+{p \over 100}\right) &|: \log \left(1+{p \over 100}\right) \\
{\log n \over \log \left(1+{p \over 100}\right)} &=& x 
\end{array}
$$
In der nach $x$ aufgelösten Beziehung kommen nur die Variablen $n$ und $p$ vor,  nicht aber das Kapital $K_0$.

\noindent {\bf Aufgabe 12. a)}
$$-2 \cdot \log_a{a} \cdot (- \log_c{c})-(-3 \cdot \log_x{x} \cdot 0)=-2 \cdot (-1) - 0 = 2$$
\noindent {\bf Aufgabe 12. b)}
$$5 \lg{a} - 4 \lg{a} - \lg{b}+ \lg{b}-\lg{a}+2,5\lg{b}=2,5\lg{b}$$

\noindent {\bf Aufgabe 13.}
Es liegt exponentielles Wachstum vor mit dem Faktor 
$$a=1+{p \over 100}= 1,05$$
Nach zehn Jahren hat sich ein Anfangskapital um den Faktor 
$$ k = 1,05^{10}=1,629 $$
vermehrt.
Der gleiche Faktor soll auch für das 7-jährige Kapital gelten; also ist der Wachstumsfaktor für ein Jahr dann:
$$
\renewcommand{\arraystretch}{2}
\begin{array}{rcll}
1,629 &=& a^7  & |\sqrt[7]{\qquad}\\
\sqrt[7]{1,629} &=& a & = 1,0722
\end{array}
$$
Hieraus errechnet sich ein Jahreszins von:
$$ p = (a-1) \cdot 100 = 7,2 \% $$

\noindent {\bf Aufgabe 14. a)}
$$
\renewcommand{\arraystretch}{2}
\begin{array}{rcll}
{5 \cdot 5^x \over 2^x } &=& 78,125 \\
5 \cdot \left({5 \over 2} \right)^x &=& 78,125 &| :5 \\
\left( {5 \over 2} \right)^x &=& 15,625 &| \log \\
x \cdot \log{5\over 2} &=& \log 15,625 & | : \log{5\over 2} \\
x &=& \frac{ \log 15,625}{\log {5\over 2}} & = 3 
\end{array}
$$
\noindent {\bf Aufgabe 14. b)}
Die Ersatzgleichung lautet:
$$
\renewcommand{\arraystretch}{2}
\begin{array}{rcll}
\sqrt{x}^{-3} &=& {1\over 27} \\
\sqrt{x}^{-3} &=& 3^{-3} & | (\qquad)^{-{1\over 3}} \\
\sqrt{x} &=& 3 &|(\qquad)^2\\
x &=& 9 
\end{array}
$$
\noindent {\bf Aufgabe 14. c)}
Die Ersatzgleichung lautet:
$$
\renewcommand{\arraystretch}{2}
\begin{array}{rcll}
x^{1\over 3} &=& 7 &|(\qquad)^3 \\
x &=& 7^{3} &= 343 
\end{array}
$$

\noindent {\bf Aufgabe 15.}\\
Der Funktionsterm für die Bakterien vom Typ A lautet (t = Stunden):
$$y_A=14000 \cdot 2^t$$
Die Bakterien vom Typ B vermehren sich um den Faktor 
$$a=1+{p\over 100} = 2,3$$
halbstündlich, also lautet der Funktionsterm: 
$$y_B=5000 \cdot 2,3^{2t}$$
Beide Terme setzen wir gleich und lösen nach t auf:
$$
\begin{array}{rcll}
14000 \cdot 2^t &=& 5000 \cdot 2,3^{2t} &| : 5000 \qquad :2^t \\
{14000 \over 5000} &=& \frac{2,3^{2t}}{2^t} \\
2,8 &=& \left( \frac{2,3^2}{2} \right)^t &| \log \\
\log 2,8 &=& t\cdot \log{2,645} &| : \log{2,645} \\
t&=& 1,06 
\end{array}
$$
Nach etwa einer Stunde sind beide Bakterienkulturen gleichauf.
 
\noindent {\bf Aufgabe 16.}

Von 2016 ($t=0$) bis 2028 sind es 12 Jahre. Schadstoffausstoß $S(t=0)=100\%$. Es ergibt sich folgende exponentielle Abnahme um jährlich $p=1,25$ \% :
$$ S(t)= 100\% \cdot (1-0,0125)^t$$
Zum Zeitpunkt ($t=12$) gilt mit eingesetzten Werten:
$$ S(12)= 100\% \cdot 0,9875^{12}= 85,9895\%$$
Insgesamt hat sich der Schadstoffausstoß dann also um $(100-86)\% = 14 \%$ verringert. Das Ziel kann auf diese Weise also erreicht werden. 

\begin{minipage}{0.5\textwidth}
\vspace{0.5 cm}
\noindent {\bf Aufgabe 17.}

a) $$5^x={1 \over 25}= 5^{-2}$$
Exponentenvergleich liefert:$$x=-2$$

b)$$ 7^c= {1\over 7^{4 \over 3}} = 7^{-{4 \over 3}}$$
Exponentenvergleich liefert: $$c=-{4 \over 3}$$

\noindent {\bf Aufgabe 18. a)}
Die Ersatzgleichung lautet:
  $$ \sqrt{a}^x=a^{2,5}$$
 $$ a^{0,5x}=a^{2,5}$$
 Daraus folgt: $$0,5x=2,5$$ und $$x=5$$.


\end{minipage}
\begin{minipage}{0.5\textwidth}
	
 \noindent {\bf Aufgabe 18. b)}
 
 Ersatzgleichung:
 $$ a^x = a^{-{1 \over 3}}$$
 $$\Rightarrow x= -{1 \over 3}$$ 
 
\noindent {\bf Aufgabe 19.}
$$x=\log_6{216}$$
$$x=3$$
\begin{center}
	\includegraphics[width=7 cm]{../../viecher/endcomic.pdf}
	
	Hier geht es zurück zum \href{https://www.okuyakl.de/math/m10logaL031/aa031.pdf}{Aufgabenblatt}
\end{center}

\end{minipage}


\end{document}

